\section{Related work}
\label{sec:related}

\paragraph{EEG foundation models and tokenizers.}
BIOT~\cite{yang2023biot} tokenizes STFT magnitude frames; LaBraM~\cite{jiang2024labram} learns discrete codes targeting per-bin amplitude and phase; CBraMod~\cite{wang2025cbramod} embeds raw 1\,s patches through a CNN with an FFT-magnitude branch and reads them with criss-cross attention; EEGPT~\cite{wang2024eegpt} and REVE~\cite{elouahidi2025reve} embed raw patches; TFM-Tokenizer~\cite{pradeepkumar2026tfm} argues that the tokenizer carries the result and transplants its codes into BIOT and LaBraM; CodeBrain~\cite{ma2025codebrain} and BandVQ~\cite{sukhbaatar2026bandvq} quantize per domain or per band. Benchmarks report that from-scratch specialists and generic time-series models tie or beat EEG-specific pretrained models~\cite{liu2026compass,kontras2026neuroatlas,xiong2025eegfmbench}, that reconstruction objectives are dominated by aperiodic low-frequency variance~\cite{kommineni2026aperiodic,yu2026fame}, and that classical features recover most of a foundation model's advantage~\cite{tang2026what}. None derives the token content from the label definitions.

\paragraph{Cross-frequency content.}
Phase--amplitude coupling has standard estimators~\cite{tort2010measuring,canolty2006high} and documented relevance in epilepsy~\cite{amiri2016phase,edakawa2016detection}; in deep models it appears as a precomputed feature or as learned amplitude--phase mixing inside an encoder~\cite{zheng2025fapex}. Cole and Voytek~\cite{cole2017brain} and Aru et al.~\cite{aru2015untangling} show that non-sinusoidal waveform shape produces apparent coupling, so a model that profits from spike-and-wave morphology profits from the structure that inflates coupling estimates. We use a coupling-derived rotation of the token only as the test transform of Section~\ref{sec:nuisance}: it is an exact isometry whose direction can be made content-independent, content-dependent, or uncorrelated with the content, which is what the hypothesis needs.

\paragraph{Factorized attention.}
Axial attention~\cite{ho2019axial} attends along one axis of a multidimensional grid at a time; CBraMod's criss-cross attention is the EEG instance over electrodes and time. The sample-complexity argument of Section~\ref{sec:factorized}, that a label which factorizes along the grid axes is learned with fewer degrees of freedom by a factorized reader, is the standard motivation for these designs applied to a grid that has a band axis.

\paragraph{Usable information and probing.}
Predictive $\mathcal{V}$-information~\cite{xu2020usable} formalizes the information a restricted predictor family can extract and shows that preprocessing can increase it although it cannot increase Shannon information; conditional probing~\cite{hewitt2021conditional} and control tasks~\cite{hewitt2019control} give the protocol for measuring it. The analytic-signal demodulation rests on Gabor's analytic signal~\cite{gabor1946theory} and Bedrosian's product theorem~\cite{bedrosian1963product}; SincNet~\cite{ravanelli2018sincnet} supplies the learned band-pass parameterization.
